\documentclass[../../main.tex]{subfiles}
\begin{document}
\begin{oframed}
  \subfile{problem}
\end{oframed}

{\bf [解答]}

\paragraph{(1)}

与えられた関数とその一階微分は
\begin{align}
  f(x) &= a\sin x + b\cos x + c\sin 2x \\
  f'(x)&=a\cos x-b\sin x+2c\cos 2x \label{eq:1}
\end{align}
である．極大値の条件から
\begin{align}
  &f'\left(\dfrac{\pi}{4}\right)=0, \quad f\left(\dfrac{\pi}{4}\right)=6\sqrt{2} \\
  \therefore
  &
  \begin{cases}
    \dfrac{\sqrt{2}}{2}(a-b)=0 \\
    \dfrac{\sqrt{2}}{2}(a+b)+c=6\sqrt{2}
  \end{cases} \label{eq:2}
\end{align}
である．次に題意の積分の条件から
\begin{align}
  5\pi
  &= \int_{0}^{2\pi}f(x)\cos x dx \\
  &= \int_{0}^{2\pi}\left[a\sin x\cos x + b\cos^2 x + c\sin 2x\cos x \right] dx \\
  &= \int_{0}^{2\pi}\left[\dfrac{a}{2}\sin 2x + \dfrac{b}{2}(1+\cos 2x) + \dfrac{c}{2}(\sin 3x+\sin x) \right] dx  \\
  &= \left[\dfrac{b}{2} \right]_{0}^{2\pi}
  = b\pi \\
  \therefore
  b = 5 \label{eq:3}
\end{align}
である．\cref{eq:2}に代入して
\begin{align}
  (a,b,c)=(5,5,\sqrt{2}) \label{eq:4}
\end{align}
となる．

\paragraph{(2)}

\cref{eq:4}を\cref{eq:1}に代入して
\begin{align}
  f(x) =  5\sin x + 5\cos x + \sqrt{2}\sin 2x
\end{align}
および
\begin{align}
  f'(x)
  &=5\cos t-5\sin t +2\sqrt{2}(\cos t^2-\sin^2 t) \\
  &=(\cos t-\sin t)\left(5+4\sin\left(x+\dfrac{\pi}{4}\right)\right)
\end{align}
だから，$f(x)$の増減表は以下のようになる．

\begin{table}[htb]
  \centering
  \caption{$f(x)$ の増減表}
  \begin{tabular}{c|ccccccc}
    $x$ & $0$ & $\cdots$ & $\dfrac{\pi}{4}$ & $\cdots$ & $\dfrac{5\pi}{4}$ & $\cdots$ & $2\pi$ \\
    \hline
    $f'(x)$ & & $+$ & $0$ & $-$ & $0$ & $+$ & \\
    \hline
    $f(x)$ & $5$ & $\nearrow$ & $6\sqrt{2}$ & $\searrow$ & $-4\sqrt{2}$ & $\nearrow$ & $5$
  \end{tabular}
\end{table}

\begin{figure}[htb]
  \centering
  \begin{tikzpicture}[x=1.3cm, y=0.45cm, >=stealth, every node/.style={font=\small}]

    % ============================================================
    % 1. 定数と座標計算
    % ============================================================
    \def\xMaxDeg{45}
    \def\xMinDeg{225}

    \pgfmathsetmacro{\piVal}{pi}
    \pgfmathsetmacro{\xMax}{\xMaxDeg * \piVal / 180}
    \pgfmathsetmacro{\xMin}{\xMinDeg * \piVal / 180}
    \pgfmathsetmacro{\xEnd}{2 * \piVal}

    \pgfmathsetmacro{\yMax}{6 * sqrt(2)}
    \pgfmathsetmacro{\yMin}{-4 * sqrt(2)}

    \coordinate (O) at (0, 0);
    \coordinate (P0) at (0, 5);
    \coordinate (Pmax) at (\xMax, \yMax);
    \coordinate (Pmin) at (\xMin, \yMin);
    \coordinate (Pend) at (\xEnd, 5);

    % ============================================================
    % 2. 座標軸
    % ============================================================
    \draw[->, thick, gray!80!black] (-0.4, 0) -- ({\xEnd + 0.6}, 0) node[right] {$x$};
    \draw[->, thick, gray!80!black] (0, {\yMin - 1.5}) -- (0, {\yMax + 1.8}) node[above] {$y$};
    \node[below left=2pt] at (O) {$O$};

    % ============================================================
    % 3. 補助点線・座標軸への投影
    % ============================================================
    % 極大点 Pmax (pi/4, 6\sqrt{2})
    \draw[densely dotted, thick, gray!70!black] (\xMax, 0) -- (Pmax);
    \draw[densely dotted, thick, gray!70!black] (0, \yMax) -- (Pmax);
    \node[below=2pt, font=\scriptsize] at (\xMax, 0) {$\frac{\pi}{4}$};
    \node[left=2pt, font=\scriptsize] at (0, \yMax) {$6\sqrt{2}$};

    % 最小点 Pmin (5pi/4, -4\sqrt{2})
    \draw[densely dotted, thick, gray!70!black] (\xMin, 0) -- (Pmin);
    \draw[densely dotted, thick, gray!70!black] (0, \yMin) -- (Pmin);
    \node[above=2pt, font=\scriptsize] at (\xMin, 0) {$\frac{5\pi}{4}$};
    \node[left=2pt, font=\scriptsize] at (0, \yMin) {$-4\sqrt{2}$};

    % 端点 x = 0, 2\pi (y = 5)
    \draw[densely dotted, thick, gray!70!black] (P0) -- (Pend) -- (\xEnd, 0);
    \node[left=2pt, font=\scriptsize] at (P0) {$5$};
    \node[below=2pt, font=\scriptsize] at (\xEnd, 0) {$2\pi$};

    % 参考点 x = \pi (目盛り)
    \draw[gray!70!black] (\piVal, 0.15) -- (\piVal, -0.15);
    \node[below=2pt, font=\scriptsize] at (\piVal, 0) {$\pi$};

    % ============================================================
    % 4. グラフ曲線 y = f(x)
    % ============================================================
    \draw[line width=1.4pt, blue!85!black, samples=120, domain=0:360, variable=\t]
    plot ({\t * \piVal / 180}, {5*sin(\t) + 5*cos(\t) + sqrt(2)*sin(2*\t)});

    \node[blue!90!black, font=\footnotesize\bfseries, above right]
    at ({\xMax + 0.4}, {\yMax - 0.7}) {$y = f(x)$};

    % ============================================================
    % 5. 点プロットとラベル
    % ============================================================
    \node[circle, fill=blue!85!black, inner sep=1.3pt] at (P0) {};
    \node[circle, fill=blue!85!black, inner sep=1.3pt] at (Pend) {};
    \node[circle, fill=red!85!black, inner sep=1.5pt] at (Pmax) {};
    \node[circle, fill=red!85!black, inner sep=1.5pt] at (Pmin) {};

    \node[above=2pt, font=\footnotesize, red!85!black] at (Pmax) {極大};
    \node[below=2pt, font=\footnotesize, red!85!black] at (Pmin) {最小};

  \end{tikzpicture}
  \caption{$y = f(x)$ のグラフ（$0 \leqq x \leqq 2\pi$）}
  \label{fig:fx_graph}
\end{figure}

従って$f(x)$は$x=\dfrac{\pi}{4}$で極大となる．これは題意の条件を満たしており十分で，
\begin{align}
  \min f(x)=f\left(\dfrac{5\pi}{4}\right)=-4\sqrt{2}
\end{align}
となる．$\cdots$(答)
\newpage
\end{document}
